\section{2018-2019学年线性代数I（H）期末答案}

\begin{enumerate}
    \item \begin{enumerate}
        \item 求核空间即求使得 $T(f(x))$ 为0矩阵的 $f(x)$ 构成的空间. 设 $f(x)=ax^{3}+bx^{2}+cx+d$，则有：
        $$
        \begin{cases}
        f(0)=d=0\\
        f(1)=a+b+c+d=0\\
        f(2)=-a+b-c+d=0
        \end{cases}
        $$
        令 $a=t$，可解得 $f(x)=t(x^{3}-x)$，故 $N(T)=\spa\{x^{3}-x\}$.

        求像空间，取 $\mathbf{R}[x]$ 的一组常用基 $1,x,x^{2},x^{3}$：
        $$
        T(1)=\begin{pmatrix}1&1\\ 1&1\end{pmatrix}, \quad T(x)=\begin{pmatrix}0&1\\ -1&0\end{pmatrix}, \quad T(x^{2})=\begin{pmatrix}0&1\\ 1&0\end{pmatrix}, \quad T(x^{3})=\begin{pmatrix}0&1\\ -1&0\end{pmatrix}
        $$
        求它们的极大线性无关组，我们发现 $T(x)=T(x^{3})$，先丢弃 $T(x^{3})$，然后令：
        $$
        k_{1}T(1)+k_{2}T(x)+k_{3}T(x^{2})=0 \implies \begin{pmatrix}k_{1}&k_{1}+k_{2}+k_{3}\\ k_{1}-k_{2}+k_{3}&k_{1}\end{pmatrix}=0
        $$
        解得 $k_1=0, k_2=0, k_3=0$. 故 $T(1),T(x),T(x^{2})$ 线性无关，从而 $R(T)=\spa\{T(1),T(x),T(x^{2})\}$.

        \item 简单验证即可，略.
    \end{enumerate}

    \item \begin{enumerate}
        \item $B$ 为上三角矩阵，则特征值为对角线元素 $2,0,1,9$. $A \sim B$ 则有相同的特征值. 则 $f(A)=A^{2}-9A+4E_{4}$ 的特征值为 $f(\lambda_{i})=\lambda_i^2 - 9\lambda_i + 4$，即为 $-10, 4, -4, 4$. 行列式的值等于特征值的乘积，为640.
        \item 由于 $A$ 有一个一重特征值0，故 $A$ 不可逆，且 $AX=0$ 的解空间维数为 $n-r(A)=1$，故 $r(A)=3, r(A^{*})=1$. 而9也为一重特征值，同理有 $r(9E-A)=3$，故 $r(A^{*})+r(9E-A)=4$.
    \end{enumerate}

    \item \begin{enumerate}
        \item 设 $\delta=(x_{1},x_{2},x_{3},x_{4})$，则由正交的定义可知：
        $$
        \begin{cases}
        \alpha\delta=x_{1}+x_{2}+x_{3}+x_{4}=0\\
        \beta\delta=-x_{1}-x_{2}+2x_{4}=0\\
        \gamma\delta=x_1-x_2=0
        \end{cases} \implies \begin{pmatrix}1&1&1&1\\ -1&-1&0&2 \\ 1&-1&0&0 \end{pmatrix}X=0
        $$
        令 $x_{1}=t$，解得 $\delta=t(1,1,-3,1)$. 又 $\delta$ 为单位向量，即 $||\delta||=1$，解得 $\delta=\pm(\dfrac{\sqrt{3}}{6},\dfrac{\sqrt{3}}{6},-\dfrac{\sqrt{3}}{2},\dfrac{\sqrt{3}}{6})$.

        \item $\|\alpha+\beta+\gamma+\delta\|=\|(1,-1,1,3) \pm (\dfrac{\sqrt{3}}{6},\dfrac{\sqrt{3}}{6},-\dfrac{\sqrt{3}}{2},\dfrac{\sqrt{3}}{6})\|=\sqrt{13}$.
    \end{enumerate}

    \item \begin{enumerate}
        \item 二次型的矩阵为 $A=\begin{pmatrix}a&0&b\\ 0&2&0\\ b&0&-2\end{pmatrix}$. 由于特征值的和等于矩阵的迹，故 $tr(A)=a=1$. 由于特征值的积等于矩阵的行列式，故 $|A|=-2(2+b^{2})=-12 \implies b=2$ (由于 $b>0$).

        \item $|\lambda E-A|=(\lambda+3)(\lambda-2)^{2}=0$，特征值为 $-3, 2, 2$.

        方程组 $(-3E-A)X = 0, (2E-A)X = 0$ 的基础解系为 $(-1,0,2)^\mathrm{T}, (0,1,0)^\mathrm{T}, (2,0,1)^\mathrm{T}$.

        正交规范化得变换矩阵 $Q=(\gamma_{1},\gamma_{2},\gamma_{3})$，其中 $\gamma_1=\dfrac{1}{\sqrt{5}}(-1,0,2)^\mathrm{T}, \gamma_2=(0,1,0)^\mathrm{T}, \gamma_3=\dfrac{1}{\sqrt{5}}(2,0,1)^\mathrm{T}$.

        \item 由于 $A$ 有负的特征值，因此非正定.
    \end{enumerate}

    \item 见\autoref{ex:非齐次解的进一步结构}. 根据该例题的结论，我们有 $AX=\beta$ 的任意解都可以表示成 $X_0, X_0+X_1, \ldots, X_0+X_s$ 的线性组合，且它们线性无关，不难证明它们的线性扩张就是所求的 $W$，基即为该向量组，维数为 $s+1$.

    \item \begin{enumerate}
        \item 取自然基 $e_1=(1,0,0), e_2=(0,1,0), e_3=(0,0,1)$，则 \[ T(e_1,e_2,e_3) = (e_1, e_2, e_3)\begin{pmatrix} 4&0&1 \\ 2&3&2 \\ 1&0&4 \end{pmatrix} \] 故特征多项式为 $(\lambda-3)^{2}(\lambda-5)$，特征值 $\lambda_{1}=\lambda_{2}=3, \lambda_{3}=5$.
        \item 对于 $\lambda_1=\lambda_2=3$，可解得特征子空间为 $\spa\{(0,1,0), (1,0,-1)\}$，是二维空间，而 $\lambda_3$ 必有一个一维的特征子空间，故几何重数等于代数重数，$T$ 可对角化.
    \end{enumerate}

    \item $r(A)=r$ 则 $AX=0$ 的解空间维数 $\dim N(A) = n-r$. 由 $r(A)+r(B)=k$ 得 $r(B)=k-r \leqslant n-r=\dim N(A)$. 要求 $AB=O$，说明 $B$ 的列向量均为 $AX=0$ 的解，那么只需要选择合适的列向量组拼接成 $B$ 即可（这一定能做到，因为 $B$ 维数不会超过解空间维数）.

    \item 见\autoref{thm:若当循环基的性质}.

    \item \begin{enumerate}
        \item 正确，因为 $A$ 的特征值不为有理数.
        \item 错误，若存在这样两组向量，则两个空间正交，从而是直和，导致矛盾.
        \item 正确，该式子结果只能为 $0,2,4$.
        \item 正确，$A$ 有特征值 $10$，对应特征向量 $(1,1,\dots,1)$，故 $f(A)$ 有特征值 $f(10)=2019$，对应每行元素之和为2019.
    \end{enumerate}
\end{enumerate}

\clearpage
